\documentclass[journal]{IEEEtran}

\usepackage{graphicx}
\usepackage{amsmath}
\usepackage{booktabs}
\usepackage{array}
\usepackage{multirow}
\usepackage{cite}

\title{Genetic-Algorithm-Optimized Mode-Mismatch Surface-Wave Suppression Using Planar Geometry Transitions}

\author{First~A.~Author,~\IEEEmembership{Student Member,~IEEE,}
Second~B.~Author,~\IEEEmembership{Member,~IEEE,}
and~Third~C.~Author,~\IEEEmembership{Senior Member,~IEEE}
\thanks{Manuscript draft prepared on March 11, 2026. Replace this block with real author affiliations, funding, and correspondence details.}
}

\begin{document}
\maketitle

\begin{abstract}
Surface waves on high-permittivity substrates degrade antenna isolation and increase unintended coupling and radiation, especially in compact printed microwave systems. Conventional suppression techniques based on electromagnetic bandgap (EBG) or defected ground structures are often constrained by wavelength-scale periodicity, multiple unit cells, narrowband resonant behavior, and increased fabrication complexity. This letter investigates a nonresonant alternative based on mode mismatch generated by abrupt transitions among several planar cross-sectional states. A genetic algorithm is used to optimize a sequence of four manufacturable geometry states with a minimum feature size of 0.3 mm. The optimization is implemented in a grounded dielectric slab model for transverse-magnetic (TM) surface waves, with fitness defined from broadband reflection. The collected TM results show a clear dependence between suppressor length and broadband reflection, with the best available case in the present dataset occurring at 6.6 mm. The optimized structures remain fully planar and do not require vias or subwavelength periodic patterning. Additional full-wave CST validation indicates that, for the same occupied length, the proposed geometry-transition method outperforms representative bandgap-based and ground-slot-based suppressors.
\end{abstract}

\begin{IEEEkeywords}
Surface waves, genetic algorithm, mode mismatch, high-permittivity substrate, broadband suppression, planar structures, MIMO decoupling.
\end{IEEEkeywords}

\section{Introduction}
High-dielectric-constant substrates are widely used to miniaturize microwave and antenna structures, but they also support stronger surface-wave excitation. These surface waves can increase mutual coupling, distort radiation patterns, reduce efficiency, and introduce unwanted parasitic radiation. In compact printed multiple-input multiple-output (MIMO) platforms, this tradeoff becomes especially severe because electrical spacing is already small.

Surface-wave suppression is commonly addressed using EBG structures, metasurfaces, or defected ground structures. Although effective in many cases, these methods usually present three practical limitations. First, their operation is often tied to wavelength-related periodic features and multiple repeated cells, which consumes valuable area and conflicts with aggressive miniaturization. Second, because many such suppressors are resonant, the achievable bandwidth is often limited. Third, implementations frequently rely on vias, dense patterning, or other fabrication steps that increase process complexity and cost.

In this work, we study a different mechanism. When a surface wave propagates from one cross-sectional geometry to another, the supported modal field changes. This mode mismatch redistributes power into reflected and radiated components without depending on a narrow resonance. The idea has appeared in decoupling-oriented structures, but prior work has generally used a limited geometry family, such as introducing slots in the ground plane of a grounded dielectric slab. Here, the geometry space is enlarged to four planar states, and a genetic algorithm (GA) is used to search for the best state sequence for broadband suppression. The resulting approach is simple, planar, and fabrication-aware, while still achieving strong suppression in both TM and, as later observed, TE cases.

\section{Structure and Optimization Method}
The baseline medium is a dielectric slab with relative permittivity $\varepsilon_r=10.2$, substrate thickness $h=1.27$ mm, and metal thickness $t=0.035$ mm. The starting configuration is a grounded dielectric slab (GDS), which is used as the reference wave-guiding structure for TM surface-wave propagation. The optimization region is partitioned into uniform longitudinal segments of 0.3 mm, chosen to remain compatible with low-cost laser-based printed-circuit-board fabrication.

Each segment can take one of four geometry states:
\begin{enumerate}
\item grounded dielectric slab with no change;
\item dielectric slab with neither top nor bottom metal in that segment;
\item dielectric slab with both top and bottom metal;
\item dielectric slab with top metal only.
\end{enumerate}

These four states can be interpreted as planar cross-sectional transitions that deliberately perturb the local supported surface-wave mode. A chromosome is therefore a one-dimensional state sequence along the propagation direction.

The TM optimization in the repository is implemented in Meep with a 2-D model, resolution 20 pixels/mm, a computational cell of $30 \times 12$ mm$^2$, and perfectly matched layers of 1 mm. Broadband excitation is applied using a Gaussian source with center frequency 0.08 and bandwidth 0.05 in Meep units. Reflection and transmission are obtained by flux monitors after normalization to the straight GDS reference. For each candidate sequence,
\begin{equation}
T(f)=\frac{P_{\mathrm{trans}}(f)}{P_{\mathrm{trans,ref}}(f)},
\qquad
R(f)=\frac{-P_{\mathrm{refl}}(f)}{P_{\mathrm{trans,ref}}(f)},
\end{equation}
and the residual term is
\begin{equation}
L(f)=1-T(f)-R(f).
\end{equation}

The current TM scripts maximize
\begin{equation}
F=\mathrm{mean}\{R(f)\}+\min\{R(f)\},
\end{equation}
which favors both high average reflection and good worst-case reflection over the operating band. The GA uses elitism, tournament selection, single-point crossover, and random mutation. The segment states at the two ends are constrained to avoid the trivial all-baseline case in the short devices.

\begin{figure}[t]
\centering
\fbox{\parbox{0.95\columnwidth}{\centering
Placeholder for concept illustration.\\
Suggested content: left-to-right propagation of a surface wave through several geometry transitions; indicate mode mismatch, partial reflection, partial radiation, and reduced transmitted surface-wave power.\\
Suggested layout/design: use a horizontal three-part schematic. Left panel: conventional high-$\varepsilon_r$ substrate with strong surface-wave propagation. Middle panel: segmented region showing states 0--3 with simple cross-sectional icons and arrows. Right panel: reflected and radiated power components highlighted with color coding; keep labels minimal and use one accent color for the proposed transition region.}}
\caption{Conceptual illustration of the proposed mode-mismatch surface-wave suppressor. Suggested layout/design: a clean horizontal schematic with one baseline section, one segmented transition region, and one output section, using arrows to distinguish incident, reflected, radiated, and transmitted power.}
\label{fig:concept}
\end{figure}

\begin{figure}[t]
\centering
\fbox{\parbox{0.95\columnwidth}{\centering
Placeholder for GA flowchart.\\
Suggested content: initialize population, decode geometry, run Meep normalization or evaluation, calculate $T/R/L$, compute fitness, rank population, selection/crossover/mutation, update best gene, and stop or continue.\\
Suggested layout/design: top-to-bottom flowchart with rounded boxes for processes and one decision diamond for the generation-stop condition; keep the normalization step visually separated as a one-time precomputation block.}}
\caption{GA optimization flow for the proposed suppressor synthesis. Suggested layout/design: a compact vertical flowchart with the one-time normalization branch on one side and the repeated evaluation loop centered to make the optimization logic easy to follow in one column.}
\label{fig:flow}
\end{figure}

\section{TM Optimization Results}
The available TM dataset spans total suppressor lengths from 0.6 to 9.0 mm in 0.3-mm increments. Table~\ref{tab:genes} summarizes the best available gene found for each tested length from the latest CSV in each folder, while Fig.~\ref{fig:length} summarizes the corresponding fitness trend. The fitness rises rapidly as the available length increases from 0.6 to about 4--5 mm, then enters a broad high-performance region. In the present dataset, the highest TM fitness is obtained at 6.6 mm, with $F=1.1616$. This indicates that strong broadband suppression can be realized with a compact transition region that is still much smaller than the wavelength-scale footprint often associated with periodic suppressors.

Representative results further illustrate the trend. At 3.0 mm, the best available sequence is
\[
[2,2,2,2,2,1,1,1,1,1],
\]
with mean reflection 0.6799 and minimum reflection 0.2918. At 6.6 mm, the best available sequence is
\[
[2,2,2,2,2,0,1,1,3,1,2,1,2,0,0,2,1,0,2,2,2,2],
\]
with mean reflection 0.7204 and minimum reflection 0.4411. At 9.0 mm, the best available sequence maintains high performance, with mean reflection 0.7307 and minimum reflection 0.4221, but does not exceed the 6.6-mm optimum in the current runs. These data suggest that longer structures provide more design freedom, but the best tradeoff is not necessarily the longest region.

Another useful observation is qualitative. The optimized sequences are not simple periodic arrangements. Instead, they mix loaded, unloaded, and asymmetric states in a nonuniform way, supporting the central premise that broadband suppression comes from cumulative mode mismatch rather than from a single narrowband resonance. This is consistent with the relatively stable reflection floor observed in the better solutions. Fig.~\ref{fig:trl} should therefore compare several representative TRL responses, for example short, medium, and near-optimal lengths, to show how the reflection floor improves as the structure gains more degrees of freedom.

\begin{table*}[t]
\centering
\caption{Best available TM gene for each tested suppressor length. Suggested layout/design: use four-column blocks across the page to save space; highlight the best overall row in bold; if the full gene strings look too dense, reduce the table font slightly or typeset the genes in fixed-width style.}
\label{tab:genes}
\begin{tabular}{c c p{2.6cm} c c p{2.6cm}}
\toprule
Length (mm) & Fitness & Best Gene & Length (mm) & Fitness & Best Gene \\
\midrule
0.6 & 0.2824 & [3, 3] & 4.8 & 0.9985 & [2,2,2,2,2,1,1,1,3,1,2,3,0,3,2,3] \\
0.9 & 0.3467 & [1,1,3] & 5.1 & 1.0874 & [2,2,0,0,2,2,2,2,2,1,1,1,0,1,0,1,1] \\
1.2 & 0.3552 & [1,1,1,2] & 5.4 & 1.1225 & [2,2,2,2,2,0,1,1,1,2,1,2,2,3,3,1,1,3] \\
1.5 & 0.4200 & [2,2,2,1,1] & 5.7 & 1.1274 & [2,2,2,2,2,0,1,1,1,1,3,0,2,2,0,1,1,1,1] \\
1.8 & 0.5006 & [2,2,2,2,0,1] & 6.0 & 1.0544 & [2,2,2,2,2,0,2,1,0,1,2,2,3,2,0,1,0,0,2,2] \\
2.1 & 0.6111 & [2,2,2,2,1,1,3] & 6.3 & 1.0736 & [2,0,2,2,2,2,0,0,3,1,1,1,1,2,1,1,1,2,2,2,2] \\
2.4 & 0.6641 & [2,2,2,2,0,3,1,3] & 6.6 & \textbf{1.1616} & [2,2,2,2,2,0,1,1,3,1,2,1,2,0,0,2,1,0,2,2,2,2] \\
2.7 & 0.7413 & [2,2,2,2,0,1,1,1,1] & 6.9 & 1.1052 & [2,2,2,2,2,0,0,1,1,0,2,2,3,3,1,1,1,1,0,1,1,2,1] \\
3.0 & 0.7631 & [2,2,2,2,2,1,1,1,1,1] & 7.2 & 1.1184 & [2,2,2,2,2,0,0,1,1,1,1,0,2,2,3,3,1,3,2,3,3,0,2,2] \\
3.3 & 0.8036 & [2,2,2,2,2,2,1,1,1,1,1] & 7.5 & 1.0991 & [2,2,2,2,2,0,2,1,1,0,2,2,2,1,3,1,1,1,3,3,3,0,2,3,3] \\
3.6 & 0.8054 & [2,2,2,2,2,2,1,1,1,1,1,1] & 7.8 & 1.0970 & [2,2,2,2,0,2,0,2,2,2,2,2,2,1,1,1,3,0,1,0,0,2,3,1,3,1] \\
3.9 & 0.9959 & [2,0,2,2,2,2,0,2,1,1,0,1,3] & 8.1 & 1.0517 & [2,2,2,2,2,0,1,0,1,0,0,0,3,1,2,0,1,1,2,0,0,0,1,0,2,0,2] \\
4.2 & 1.0226 & [2,0,2,2,2,2,0,2,1,1,1,1,1,1] & 8.4 & 1.0899 & [2,2,2,2,2,0,1,1,1,2,2,3,2,0,3,3,0,2,3,2,2,0,1,0,0,3,0,2] \\
4.5 & 0.8689 & [2,2,2,0,2,2,2,2,0,1,1,1,1,1,1] & 8.7 & 1.0930 & [2,2,2,2,2,0,1,1,3,0,1,3,1,3,1,2,1,2,3,3,3,2,2,3,0,1,1,3,2] \\
\multicolumn{3}{c}{} & 9.0 & 1.1528 & [2,0,0,2,2,2,2,2,0,0,1,0,1,0,2,2,3,2,1,1,0,3,3,2,0,0,0,3,3,1] \\
\bottomrule
\end{tabular}
\end{table*}

\begin{figure}[t]
\centering
\includegraphics[width=\columnwidth]{../GA_SW_suppression_TM/fitness_over_length_TM.png}
\caption{Best available TM fitness versus suppressor length, extracted from the latest CSV found in each length folder. The present dataset peaks at 6.6 mm. Suggested layout/design: keep this as a clean single-column line plot with markers, label only a few representative genes to avoid clutter, and optionally annotate the 6.6-mm optimum with an arrow callout.}
\label{fig:length}
\end{figure}

\begin{figure}[t]
\centering
\fbox{\parbox{0.95\columnwidth}{\centering
Placeholder for representative TRL curves.\\
Suggested content: three to four representative designs, such as 3.0 mm, 4.2 mm, 6.6 mm, and 9.0 mm, plotted over the same frequency range.\\
Suggested layout/design: use a $2\times2$ panel figure if you want one design per panel, or overlay only the $R$ curves in one panel and place corresponding $T$ and $L$ in smaller insets. For readability in AWPL, avoid plotting too many full TRL triplets in a single axis.}}
\caption{Representative transmission, reflection, and loss responses for selected optimized genes. Suggested layout/design: either a $2\times2$ panel figure with one design per panel or a compact overlay emphasizing $R$ first, since reflection is the optimization target and the key story.}
\label{fig:trl}
\end{figure}

\section{Discussion and Comparison}
Compared with periodic EBG-based suppressors, the proposed method has several advantages. First, it does not require several repeated wavelength-related cells, which is attractive for tightly packed PCB antennas and compact decoupling regions. Second, the mechanism is nonresonant in the usual periodic-bandgap sense, so broadband suppression is easier to realize. Third, the geometry can remain fully planar, without vias and without dense subwavelength perforations, making fabrication simpler and potentially cheaper.

This geometry-transition viewpoint also provides a useful bridge between surface-wave control and antenna decoupling. In many compact MIMO designs, decoupling structures succeed because they interrupt the preferred current or surface-wave path. The present work formalizes that idea as a mode-mismatch optimization problem over a larger geometry space. Although the GA optimization was primarily carried out for TM excitation, the same optimized or closely related geometries were later found to provide strong suppression for TE surface waves as well, which supports the broadband and polarization-robust nature of the concept. Quantitative TE plots can be added in the final version if desired.

Full-wave CST simulations further show that, for the same occupied length, the proposed suppressor achieves stronger suppression than representative bandgap-based and ground-slot-based methods. Because those CST comparison figures are not yet stored in this repository, two validation figures are reserved below: one for the direct suppression comparison and one for a field or current-distribution comparison that visually explains why the proposed structure works. A concise comparison table can also be included if page space permits.

\begin{figure}[t]
\centering
\fbox{\parbox{0.95\columnwidth}{\centering
Placeholder for first CST validation figure.\\
Suggested content: proposed method vs. EBG-based suppressor vs. ground-slot-based suppressor at equal occupied length, using $|S_{21}|$, transmission loss, or normalized power transmission over frequency.\\
Suggested layout/design: use one shared axis with three clear line styles and a shaded band for the main operating region, or two stacked panels if both TM and TE are shown.}}
\caption{CST comparison of the proposed suppressor against reference suppression methods at the same occupied length. Suggested layout/design: one main comparison plot with consistent colors and line styles, and add a small inset listing each structure's total length and fabrication complexity.}
\label{fig:cst_compare}
\end{figure}

\begin{figure}[t]
\centering
\fbox{\parbox{0.95\columnwidth}{\centering
Placeholder for second CST validation figure.\\
Suggested content: electric-field, magnetic-field, or surface-current distribution for the proposed structure and one reference structure, showing stronger attenuation or reflection inside the optimized transition region.\\
Suggested layout/design: use side-by-side field maps with identical color scales and add a dashed box around the optimized region; if space is tight, use one field map plus one 1-D extracted power-decay curve.}}
\caption{CST field-distribution validation of the suppression mechanism. Suggested layout/design: side-by-side field maps with the same color bar so the viewer can directly compare how the proposed transition region blocks or reflects the surface wave relative to a reference design.}
\label{fig:cst_field}
\end{figure}

\section{Conclusion}
This letter presented a GA-based, mode-mismatch surface-wave suppressor formed by longitudinal transitions among four planar geometry states. Unlike conventional periodic or resonant suppressors, the proposed structure is compact, broadband, and fabrication-friendly. The TM optimization results show a strong length dependence and identify a 6.6-mm design as the best case in the current dataset. The optimized sequences are nonperiodic, indicating that cumulative modal disruption is the dominant suppression mechanism. Preliminary TE behavior and CST comparisons further suggest that the approach is broadly effective and competitive with more fabrication-intensive alternatives. This makes the method promising for miniaturized high-permittivity microwave and MIMO platforms.

\section*{Notes for Final Submission}
\begin{itemize}
\item If your final objective was changed to $F=\mathrm{mean}(R)+0.5\min(R)$ in later runs, update (3) and the corresponding discussion.
\item Replace the placeholder CST figure with your actual comparison results.
\item Add the final TE validation figure or a short sentence with numerical values if you want the dual-polarization claim to be stronger.
\item Replace the draft references below with the exact papers you want to cite.
\end{itemize}

\begin{thebibliography}{99}
\bibitem{ref1}
F.~Yang and Y.~Rahmat-Samii, ``Microstrip antennas integrated with electromagnetic band-gap (EBG) structures: a low mutual coupling design for array applications,'' \emph{IEEE Trans. Antennas Propag.}, placeholder entry for final editing.

\bibitem{ref2}
J.~G.~Yook, H.~Lee, and H.~Kim, ``Enhanced patch-antenna performance by suppressing surface waves using photonic-bandgap substrates,'' \emph{IEEE Microw. Guided Wave Lett.}, placeholder entry for final editing.

\bibitem{ref3}
``Mutual coupling reduction in MIMO antenna system using EBG structures,'' \emph{IEEE conference paper}, IEEE Xplore document 6290217, placeholder entry for final editing. Available: \texttt{https://ieeexplore.ieee.org/document/6290217/}

\bibitem{ref4}
``Mutual Coupling Suppression With Decoupling Ground for Massive MIMO Antenna Arrays,'' \emph{IEEE Access}, IEEE Xplore document 8737717, placeholder entry for final editing. Available: \texttt{https://ieeexplore.ieee.org/document/8737717/}

\bibitem{ref5}
``Defected ground structure with two resonances for decoupling of dual-band MIMO antenna,'' \emph{IEEE conference paper}, IEEE Xplore document 8072865, placeholder entry for final editing. Available: \texttt{https://ieeexplore.ieee.org/document/8072865/}

\bibitem{ref6}
``Mutual Coupling Reduction of a Two-element MIMO Antenna System Using Defected Ground Structure,'' \emph{IEEE conference paper}, IEEE Xplore document 9330404, placeholder entry for final editing. Available: \texttt{https://ieeexplore.ieee.org/document/9330404/}
\end{thebibliography}

\end{document}
